\section{Background and Motivation}
\label{sec:background}

\paragraph{Mobile agents act through the UI.} Agents such as AppAgent, AutoDroid, and the Mobile-Agent family drive Android through the accessibility tree, screenshots, and synthesized gestures \cite{zhang2023appagent,wen2023autodroid,wang2024mobileagent,ye2025mobileagentv3}. Everything the agent perceives is therefore an observation channel, and everything it does is an actuation channel. A secret that must be entered crosses both.

\paragraph{The motivating attack.} In the OTP-forwarding pattern, the agent encounters a message or notification containing a verification code together with an injected instruction to forward it. The agent reads the code from the rendered text and sends it to an attacker-controlled endpoint \cite{du2025mobileattacks,greshake2023indirect}. We reproduced the core mechanism on an Android emulator with an open notification-listener agent: the code was read from the notification and POSTed to an isolated endpoint (Section~\ref{sec:evidence}). This establishes that the attack surface is reachable with stock, non-malicious-looking components.

\paragraph{Why existing defenses fail.} \emph{Placeholder anonymization} replaces rendered sensitive text with a local token and resolves it at \texttt{type} time. It removes the literal value from the current observation but depends on detector recall, is session-scoped, and does not constrain \emph{which} field receives the resolved value or \emph{which} origin may request it. An injected instruction can redirect the resolution \cite{zhao2026anonymization}. \emph{Platform autofill and Credential Manager} are app-initiated: the consuming app must drive them, and they trust that app, so an accessibility agent cannot invoke them and a hostile field is not distinguished \cite{androidAutofill,androidAutofillServices,androidCredentialManager}. \emph{Server-side agentic-payment tokens} cover payments only and live outside the OS, leaving OTP, phone, email, address, and clipboard unprotected.

\paragraph{Design forces.} A defense that bends the attack curve must (1) mint the secret at the authoritative source, so security does not depend on detector recall; (2) bind it to the destination and task that may consume it, so an injection cannot reroute it; (3) give it a typed lifecycle, so single-use secrets cannot be replayed; and (4) redact the resolved value afterwards, because the entered value reappears in the next frame.
